% paper/sections/results.tex
% Discussion / analysis: cross-regime findings, Transformer variant, limitations

\section{Discussion}
\label{sec:discussion}

\paragraph{Why the learned policy dominates at small $k$.}
At $k{=}1$ (budget $B_0{=}0.30$) the entire session is defined by one critical
decision: find a buyer who will accept at a price that both replenishes the budget
\emph{and} yields high revenue.
Greedy+Budget selects by degree regardless of pricing; if the first offer is
rejected the session ends bankrupt.
The learned policy uses the full episode context — valuation estimates, influence
state, budget ratio — to set a discount calibrated to ensure acceptance, achieving
$305.6$ revenue vs.\ $7.4$ for Greedy+Budget. %SOURCE: unified_sweep.json / dp_upgrade_eval.json
Cal-DP v3 also exploits this structure ($10.6$) but cannot dynamically adapt
pricing to individual buyer responses, capping its advantage.

\paragraph{Why Cal-DP leads at mid-$k$ on dense graphs.}
On the Rice-Facebook network ($|\calE|{=}4{,}752$ vs.\ $2{,}377$ for FF $n{=}1000$)
the influence model assigns high valuations to many nodes, not just high-degree
hubs.
Cal-DP v2 exploits a globally computed value function over the budget grid,
effectively finding near-optimal buyer orderings under the dense influence
structure.
The learned policy, trained on sparser Forest-Fire graphs, does not adapt its
pricing calibration to the denser topology at mid-$k$ (capital-transitional
regime).
This failure mode is an honest finding and motivates future transfer or
fine-tuning experiments.

\paragraph{Large-$k$ specialist and the slack regime.}
The seed-matched evaluation confirms that the gap between the unified learned
policy ($407.0$ at $k{=}40$) and Cal-DP ($448.0$) is a training-range artifact,
not a fundamental limitation.
Restricting training to $k \in [16, 40]$ and warm-starting from the unified
model eliminates the opposing-gradient tension, yielding $473.1\,{\pm}\,3.5$
at $k{=}40$ — exceeding both Cal-DP ($448.0$) and Greedy+Budget ($448.7$).
The specialist is deployed for $k{\ge}16$; the unified model handles $k{\le}15$,
where Cal-DP retains the strongest position in the transitional band.

\paragraph{Transformer variant.}
We also trained a Graph Transformer variant (Rev-GNN-Transformer) replacing the
SAGEConv layers with TransformerConv~\citep{shi2021masked}.
On FF $n{=}1000$ (unconstrained) it achieves $463.8 \pm 5.3$, %SOURCE: transformer_gate_a.json
narrowly exceeding Rev-GNN-LSTM ($462.6$, Gate~A PASS). %SOURCE: CLAUDE.md
In the budget-constrained setting, the Transformer wins $5/10$ $k$-values on FF
and $2/10$ on Rice-Facebook (Gate~B PASS). %SOURCE: gate_b_eval_v2.json
Given the marginal advantage and higher inference cost, we report
Rev-GNN-Transformer as a supplemental result; all primary claims use Rev-GNN-LSTM.

\paragraph{Limitations.}
\begin{enumerate}[leftmargin=*,topsep=2pt,itemsep=1pt]
  \item \emph{Single influence model}: results are for the Rayleigh valuation model
    of \citet{babaei2013revenue}; other valuation functions (e.g., linear threshold,
    independent cascade-based) are not tested.
  \item \emph{Small training graphs}: the model is trained on $n \le 440$ and
    evaluated up to $n{=}2000$; very large real networks ($n > 10^4$) are not
    evaluated.
  \item \emph{Budget model}: the replenishable capital model is a simplification;
    real firms face compound costs (shipping, returns, marketing overhead) not
    captured here.
  \item \emph{Edge weight uncertainty}: we use 200 Monte Carlo samples to estimate
    valuations; the ground-truth weights are assumed unknown, which may not hold
    in practice.
\end{enumerate}
